\documentclass[12pt,a4paper]{ctexbook}
\usepackage{geometry}
\geometry{margin=2.2cm}
\usepackage{setspace}
\setstretch{1.35}
\usepackage{hyperref}
\title{全唐诗选·poet.tang.0}
\author{古典文献整理}
\date{}
\begin{document}
\maketitle
\tableofcontents
\newpage
\section{帝京篇十首 一}
\textbf{太宗皇帝}\par\vspace{0.5em}
秦川雄帝宅，函谷壯皇居。\par
綺殿千尋起，離宮百雉餘。\par
連甍遙接漢，飛觀迥凌虛。\par
雲日隱層闕，風煙出綺疎。\par

\section{帝京篇十首 二}
\textbf{太宗皇帝}\par\vspace{0.5em}
巖廊罷機務，崇文聊駐輦。\par
玉匣啓龍圖，金繩披鳳篆。\par
韋編斷仍續，縹帙舒還卷。\par
對此乃淹留，欹案觀墳典。\par

\section{帝京篇十首 三}
\textbf{太宗皇帝}\par\vspace{0.5em}
移步出詞林，停輿欣武宴。\par
琱弓寫明月，駿馬疑流電。\par
驚雁落虛弦，啼猿悲急箭。\par
閱賞誠多美，於茲乃忘倦。\par

\section{帝京篇十首 四}
\textbf{太宗皇帝}\par\vspace{0.5em}
鳴笳臨樂館，眺聽歡芳節。\par
急管韻朱絃，清歌凝白雪。\par
彩鳳肅來儀，玄鶴紛成列。\par
去茲鄭衛聲，雅音方可悅。\par

\section{帝京篇十首 五}
\textbf{太宗皇帝}\par\vspace{0.5em}
芳辰追逸趣，禁苑信多奇。\par
橋形通漢上，峰勢接雲危。\par
煙霞交隱映，花鳥自參差。\par
何如肆轍跡？萬里賞瑤池。\par

\section{帝京篇十首 六}
\textbf{太宗皇帝}\par\vspace{0.5em}
飛蓋去芳園，蘭橈遊翠渚。\par
萍間日彩亂，荷處香風舉。\par
桂楫滿中川，弦歌振長嶼。\par
豈必汾河曲，方爲歡宴所。\par

\section{帝京篇十首 七}
\textbf{太宗皇帝}\par\vspace{0.5em}
落日雙闕昏，回輿九重暮。\par
長煙散初碧，皎月澄輕素。\par
搴幌翫琴書，開軒引雲霧。\par
斜漢耿層閣，清風搖玉樹。\par

\section{帝京篇十首 八}
\textbf{太宗皇帝}\par\vspace{0.5em}
歡樂難再逢，芳辰良可惜。\par
玉酒泛雲罍，蘭殽陳綺席。\par
千鍾合堯禹，百獸諧金石。\par
得志重寸陰，忘懷輕尺璧。\par

\section{帝京篇十首 九}
\textbf{太宗皇帝}\par\vspace{0.5em}
建章歡賞夕，二八盡妖妍。\par
羅綺昭陽殿，芬芳玳瑁筵。\par
珮移星正動，扇掩月初圓。\par
無勞上懸圃，即此對神仙。\par

\section{帝京篇十首 十}
\textbf{太宗皇帝}\par\vspace{0.5em}
以茲遊觀極，悠然獨長想。\par
披卷覽前蹤，撫躬尋既往。\par
望古茅茨約，瞻今蘭殿廣。\par
人道惡高危，虛心戒盈蕩。\par
奉天竭誠敬，臨民思惠養。\par
納善察忠諫，明科慎刑賞。\par
六五誠難繼，四三非易仰。\par
廣待淳化敷，方嗣云亭響。\par

\section{飲馬長城窟行}
\textbf{太宗皇帝}\par\vspace{0.5em}
塞外悲風切，交河冰已結。\par
瀚海百重波，陰山千里雪。\par
迥戍危烽火，層巒引高節。\par
悠悠卷斾旌，飲馬出長城。\par
寒沙連騎跡，朔吹斷邊聲。\par
胡塵清玉塞，羌笛韻金鉦。\par
絕漠干戈戢，車徒振原隰。\par
都尉反龍堆，將軍旋馬邑。\par
揚麾氛霧靜，紀石功名立。\par
荒裔一戎衣，靈臺凱歌入。\par

\section{執契靜三邊}
\textbf{太宗皇帝}\par\vspace{0.5em}
執契靜三邊，持衡臨萬姓。\par
玉彩輝關燭，金華流日鏡。\par
無爲宇宙清，有美璇璣正。\par
皎佩星連景，飄衣雲結慶。\par
戢武耀七德，昇文輝九功。\par
煙波澄舊碧，塵火息前紅。\par
霜野韜蓮劒，關城罷月弓。\par
錢綴榆天合，新城柳塞空。\par
花銷蔥嶺雪，縠盡流沙霧。\par
秋駕轉兢懷，春冰彌軫慮。\par
書絕龍庭羽，烽休鳳穴戍。\par
衣宵寢二難，食旰餐三懼。\par
翦暴興先廢，除兇存昔亡。\par
圓蓋歸天壤，方輿入地荒。\par
孔海池京邑，雙河沼帝鄉。\par
循躬思勵己，撫俗媿時康。\par
元首佇鹽梅，股肱惟輔弼。\par
羽賢崆嶺四，翼聖襄城七。\par
澆俗庶反淳，替文聊就質。\par
已知隆至道，共歡區宇一。\par

\section{正日臨朝}
\textbf{太宗皇帝}\par\vspace{0.5em}
條風開獻節，灰律動初陽。\par
百蠻奉遐賮，萬國朝未央。\par
雖無舜禹迹，幸欣天地康。\par
車軌同八表，書文混四方。\par
赫奕儼冠蓋，紛綸盛服章。\par
羽旄飛馳道，鐘鼓震巖廊。\par
組練輝霞色，霜戟耀朝光。\par
晨宵懷至理，終媿撫遐荒。\par

\section{幸武功慶善宮}
\textbf{太宗皇帝}\par\vspace{0.5em}
壽丘惟舊跡，酆邑乃前基。\par
粵予承累聖，懸弧亦在茲。\par
弱齡逢運改，提劒鬱匡時。\par
指麾八荒定，懷柔萬國夷。\par
梯山咸入款，駕海亦來思。\par
單于陪武帳，日逐衛文㮰。\par
端扆朝四岳，無爲任百司。\par
霜節明秋景，輕冰結水湄。\par
芸黃徧原隰，禾穎積京畿。\par
共樂還鄉宴，歡比大風詩。\par

\section{重幸武功}
\textbf{太宗皇帝}\par\vspace{0.5em}
代馬依朔吹，驚禽愁昔叢。\par
況茲承眷德，懷舊感深衷。\par
積善忻餘慶，暢武悅成功。\par
垂衣天下治，端拱車書同。\par
白水巡前跡，丹陵幸舊宮。\par
列筵歡故老，高宴聚新豐。\par
駐蹕撫田畯，回輿訪牧童。\par
瑞氣縈丹闕，祥煙散碧空。\par
孤嶼含霜白，遙山帶日紅。\par
於焉歡擊筑，聊以詠南風。\par

\section{經破薛舉戰地}
\textbf{太宗皇帝}\par\vspace{0.5em}
昔年懷壯氣，提戈初仗節。\par
心隨朗日高，志與秋霜潔。\par
移鋒驚電起，轉戰長河決。\par
營碎落星沈，陣卷橫雲裂。\par
一揮氛沴靜，再舉鯨鯢滅。\par
於茲俯舊原，屬目駐華軒。\par
沈沙無故迹，減竈有殘痕。\par
浪霞穿水淨，峰霧抱蓮昏。\par
世途亟流易，人事殊今昔。\par
長想眺前蹤，撫躬聊自適。\par

\section{過舊宅二首 一}
\textbf{太宗皇帝}\par\vspace{0.5em}
新豐停翠輦，譙邑駐鳴笳。\par
園荒一徑斷，苔古半階斜。\par
前池消舊水，昔樹發今花。\par
一朝辭此地，四海遂爲家。\par

\section{過舊宅二首 二}
\textbf{太宗皇帝}\par\vspace{0.5em}
金輿巡白水，玉輦駐新豐。\par
紐落藤披架，花殘菊破叢。\par
葉鋪荒草蔓，流竭半池空。\par
紉珮蘭凋徑，舒圭葉翦桐。\par
昔地一蕃內，今宅九圍中。\par
架海波澄鏡，韜戈器反農。\par
八表文同軌，無勞歌大風。\par

\section{還陝述懷}
\textbf{太宗皇帝}\par\vspace{0.5em}
慨然撫長劒，濟世豈邀名。\par
星旂紛電舉，日羽肅天行。\par
徧野屯萬騎，臨原駐五營。\par
登山麾武節，背水縱神兵。\par
在昔戎戈動，今來宇宙平。\par

\section{入潼關}
\textbf{太宗皇帝}\par\vspace{0.5em}
崤函稱地險，襟帶壯兩京。\par
霜峰直臨道，冰河曲繞城。\par
古木參差影，寒猿斷續聲。\par
冠蓋往來合，風塵朝夕驚。\par
高談先馬度，僞曉預雞鳴。\par
棄繻懷遠志，封泥負壯情。\par
別有真人氣，安知名不名。\par

\section{於北平作}
\textbf{太宗皇帝}\par\vspace{0.5em}
翠野駐戎軒，盧龍轉征斾。\par
遙山麗如綺，長流縈似帶。\par
海氣百重樓，巖松千丈蓋。\par

\section{遼城望月}
\textbf{太宗皇帝}\par\vspace{0.5em}
玄兔月初明，澄輝照遼碣。\par
映雲光暫隱，隔樹花如綴。\par
魄滿桂枝圓，輪虧鏡彩缺。\par
臨城却影散，帶暈重圍結。\par
駐蹕俯九都，停觀妖氛滅。\par

\section{春日登陝州城樓俯眺原野迴丹碧綴煙霞密翠斑紅芳菲花柳即目川岫聊以命篇}
\textbf{太宗皇帝}\par\vspace{0.5em}
碧原開霧隰，綺嶺峻霞城。\par
煙峰高下翠，日浪淺深明。\par
斑紅粧橤樹，圓青壓溜荆。\par
迹巖勞傅想，窺野訪莘情。\par
巨川何以濟，舟楫佇時英。\par

\section{春日玄武門宴羣臣}
\textbf{太宗皇帝}\par\vspace{0.5em}
韶光開令序，淑氣動芳年。\par
駐輦華林側，高宴柏梁前。\par
紫庭文珮滿，丹墀衮紱連。\par
九夷簉瑤席，五狄列瓊筵。\par
娛賓歌湛露，廣樂奏鈞天。\par
清尊浮綠醑，雅曲韻朱弦。\par
粵余君萬國，還慙撫八埏。\par
庶幾保貞固，虛己厲求賢。\par

\section{登三臺言志}
\textbf{太宗皇帝}\par\vspace{0.5em}
未央初壯漢，阿房昔侈秦。\par
在危猶騁麗，居奢遂役人。\par
豈如家四海，日宇罄朝倫。\par
扇天裁戶舊，砌地翦基新。\par
引月擎宵桂，飄雲逼曙鱗。\par
露除光炫玉，霜闕映雕銀。\par
舞接花梁燕，歌迎鳥路塵。\par
鏡池波太液，莊苑麗宜春。\par
作異甘泉日，停非路寢辰。\par
念勞慙逸己，居曠返勞神。\par
所欣成大厦，宏材佇渭濱。\par

\section{出獵}
\textbf{太宗皇帝}\par\vspace{0.5em}
楚王雲夢澤，漢帝長楊宮。\par
豈若因農暇，閱武出轘嵩。\par
三驅陳銳卒，七萃列材雄。\par
寒野霜氛白，平原燒火紅。\par
琱戈夏服箭，羽騎綠沈弓。\par
怖獸潛幽壑，驚禽散翠空。\par
長煙晦落景，灌木振嚴風。\par
所爲除民瘼，非是悅林叢。\par

\section{冬狩}
\textbf{太宗皇帝}\par\vspace{0.5em}
烈烈寒風起，慘慘飛雲浮。\par
霜濃凝廣隰，冰厚結清流。\par
金鞍移上苑，玉勒騁平疇。\par
旌旗四望合，罝羅一面求。\par
楚踣爭兕殪，秦亡角鹿愁。\par
獸忙投密樹，鴻驚起礫洲。\par
騎斂原塵靜，戈迴嶺日收。\par
心非洛汭逸，意在渭濱游。\par
禽荒非所樂，撫轡更招憂。\par

\section{春日望海}
\textbf{太宗皇帝}\par\vspace{0.5em}
披襟眺滄海，憑軾翫春芳。\par
積流橫地紀，疎派引天潢。\par
仙氣凝三嶺，和風扇八荒。\par
拂潮雲布色，穿浪日舒光。\par
照岸花分彩，迷雲雁斷行。\par
懷卑運深廣，持滿守靈長。\par
有形非易測，無源詎可量。\par
洪濤經變野，翠島屢成桑。\par
之罘思漢帝，碣石想秦皇。\par
霓裳非本意，端拱且圖王。\par

\section{臨洛水}
\textbf{太宗皇帝}\par\vspace{0.5em}
春蒐馳駿骨，總轡俯長河。\par
霞處流縈錦，風前瀁卷羅。\par
水花翻照樹，堤蘭倒插波。\par
豈必汾陰曲，秋雲發棹歌。\par

\section{望終南山}
\textbf{太宗皇帝}\par\vspace{0.5em}
重巒俯渭水，碧嶂插遙天。\par
出紅扶嶺日，入翠貯巖煙。\par
疊松朝若夜，複岫闕疑全。\par
對此恬千慮，無勞訪九仙。\par

\section{元日}
\textbf{太宗皇帝}\par\vspace{0.5em}
高軒曖春色，邃閣媚朝光。\par
彤庭飛綵斾，翠幌曜明璫。\par
恭己臨四極，垂衣馭八荒。\par
霜戟列丹陛，絲竹韻長廊。\par
穆矣熏風茂，康哉帝道昌。\par
繼文遵後軌，循古鑒前王。\par
草秀故春色，梅豔昔年粧。\par
巨川思欲濟，終以寄舟航。\par

\section{初春登樓即目觀作述懷}
\textbf{太宗皇帝}\par\vspace{0.5em}
憑軒俯蘭閣，眺矚散靈襟。\par
綺峰含翠霧，照日蕊紅林。\par
鏤丹霞錦岫，殘素雪斑岑。\par
拂浪堤垂柳，嬌花鳥續吟。\par
連甍豈一拱，衆幹如千尋。\par
明非獨材力，終藉棟梁深。\par
彌懷矜樂志，更懼戒盈心。\par
媿制勞居逸，方規十產金。\par

\section{首春}
\textbf{太宗皇帝}\par\vspace{0.5em}
寒隨窮律變，春逐鳥聲開。\par
初風飄帶柳，晚雪間花梅。\par
碧林青舊竹，綠沼翠新苔。\par
芝田初雁去，綺樹巧鸎來。\par

\section{初晴落景}
\textbf{太宗皇帝}\par\vspace{0.5em}
晚霞聊自怡，初晴彌可喜。\par
日晃百花色，風動千林翠。\par
池魚躍不同，園鳥聲還異。\par
寄言博通者，知予物外志。\par

\section{初夏}
\textbf{太宗皇帝}\par\vspace{0.5em}
一朝春夏改，隔夜鳥花遷。\par
陰陽深淺葉，曉夕重輕煙。\par
哢鸎猶響殿，橫絲正網天。\par
珮高蘭影接，綬細草紋連。\par
碧鱗驚櫂側，玄燕舞檐前。\par
何必汾陽處，始復有山泉。\par

\section{度秋}
\textbf{太宗皇帝}\par\vspace{0.5em}
夏律昨留灰，秋箭今移晷。\par
峨嵋岫初出，洞庭波漸起。\par
桂白發幽巖，菊黃開灞涘。\par
運流方可歎，含毫屬微理。\par

\section{儀鸞殿早秋}
\textbf{太宗皇帝}\par\vspace{0.5em}
寒驚薊門葉，秋發小山枝。\par
松陰背日轉，竹影避風移。\par
提壺菊花岸，高興芙蓉池。\par
欲知涼氣早，巢空燕不窺。\par

\section{秋日即目}
\textbf{太宗皇帝}\par\vspace{0.5em}
爽氣浮丹闕，秋光澹紫宮。\par
衣碎荷疎影，花明菊點叢。\par
袍輕低草露，蓋側舞松風。\par
散岫飄雲葉，迷路飛煙鴻。\par
砌冷蘭凋佩，閨寒樹隕桐。\par
別鶴棲琴裏，離猿啼峽中。\par
落野飛星箭，弦虛半月弓。\par
芳菲夕霧起，暮色滿房櫳。\par

\section{山閣晚秋}
\textbf{太宗皇帝}\par\vspace{0.5em}
山亭秋色滿，巖牖涼風度。\par
疎蘭尚染煙，殘菊猶承露。\par
古石衣新苔，新巢封古樹。\par
歷覽情無極，咫尺輸光暮。\par

\section{秋暮言志}
\textbf{太宗皇帝}\par\vspace{0.5em}
朝光浮燒野，霜華淨碧空。\par
結浪冰初鏡，在逕菊方叢。\par
約嶺煙深翠，分旗霞散紅。\par
抽思滋泉側，飛想傅巖中。\par

\section{喜雪}
\textbf{太宗皇帝}\par\vspace{0.5em}
碧昏朝合霧，丹卷暝韜霞。\par
結葉繁雲色，凝璚徧雪華。\par
光樓皎若粉，映幕集疑沙。\par
泛柳飛飛絮，粧梅片片花。\par
照璧臺圓月，飄珠箔穿露。\par
瑤潔短長階，玉藂高下樹。\par
映桐珪累白，縈峰蓮抱素。\par
斷續氣將沈，徘徊歲云暮。\par
懷珍媿隱德，表瑞佇豐年。\par
橤間飛禁苑，鶴處舞伊川。\par
儻詠幽蘭曲，同懽黃竹篇。\par

\section{秋日斅庾信體}
\textbf{太宗皇帝}\par\vspace{0.5em}
嶺銜宵月桂，珠穿曉露叢。\par
蟬啼覺樹冷，螢火不溫風。\par
花生圓菊橤，荷盡戲魚通。\par
晨浦鳴飛雁，夕渚集栖鴻。\par
颯颯高天吹，氛澄下熾空。\par

\section{賦尚書}
\textbf{太宗皇帝}\par\vspace{0.5em}
崇文時駐步，東觀還停輦。\par
輟膳翫三墳，暉燈披五典。\par
寒心覩肉林，飛魄看沈湎。\par
縱情昏主多，克己明君鮮。\par
滅身資累惡，成名由積善。\par
既承百王末，戰兢隨歲轉。\par

\section{詠司馬彪續漢志}
\textbf{太宗皇帝}\par\vspace{0.5em}
二儀初創象，三才乃分位。\par
非惟樹司牧，固亦垂文字。\par
緜代更膺期，芳圖無輟記。\par
炎漢承君道，英謨纂神器。\par
潛龍既可躍，逵兔奚難致。\par
前史殫妙詞，後昆沈雅思。\par
書言揚盛跡，補闕興洪志。\par
川谷猶舊途，郡國開新意。\par
梅山未覺朽，穀水誰云異。\par
車服隨名表，文物因時置。\par
鳳戟翼康衢，鑾輿總柔轡。\par
清濁必能澄，洪纖幸無棄。\par
觀儀不失序，遵禮方由事。\par
政宣竹律和，時平玉條備。\par
文囿雕奇彩，藝門蘊深致。\par
雲飛星共流，風揚月兼至。\par
類禋遵令典，壇壝資良地。\par
五勝竟無違，百司誠有庇。\par
粵予承暇景，談叢引泉祕。\par
討論窮義府，看覈披經笥。\par
大辨良難仰，小學終先匱。\par
聞道諒知榮，含毫孰忘媿。\par

\section{詠風}
\textbf{太宗皇帝}\par\vspace{0.5em}
蕭條起關塞，搖颺下蓬瀛。\par
拂林花亂彩，響谷鳥分聲。\par
披雲羅影散，泛水織文生。\par
勞歌大風曲，威加四海清。\par

\section{詠雨}
\textbf{太宗皇帝}\par\vspace{0.5em}
罩雲飄遠岫，噴雨泛長河。\par
低飛昏嶺腹，斜足灑巖阿。\par
泫叢珠締葉，起溜鏡圖波。\par
濛柳添絲密，含吹織空羅。\par

\section{詠雪}
\textbf{太宗皇帝}\par\vspace{0.5em}
潔野凝晨曜，裝墀帶夕暉。\par
集條分樹玉，拂浪影泉璣。\par
色灑妝臺粉，花飄綺席衣。\par
入扇縈離匣，點素皎殘機。\par

\section{賦得夏首啓節}
\textbf{太宗皇帝}\par\vspace{0.5em}
北闕三春晚，南榮九夏初。\par
黃鶯弄漸變，翠林花落餘。\par
瀑流還響谷，猿啼自應虛。\par
早荷向心卷，長楊就影舒。\par
此時歡不極，調軫坐相於。\par

\section{賦得白日半西山}
\textbf{太宗皇帝}\par\vspace{0.5em}
紅輪不暫駐，烏飛豈復停。\par
岑霞漸漸落，溪陰寸寸生。\par
藿葉隨光轉，葵心逐照傾。\par
晚煙含樹色，棲鳥雜流聲。\par

\section{置酒坐飛閣}
\textbf{太宗皇帝}\par\vspace{0.5em}
高軒臨碧渚，飛檐迥架空。\par
餘花攢鏤檻，殘柳散雕櫳。\par
岸菊初含橤，園梨始帶紅。\par
莫慮崑山暗，還共盡杯中。\par

\section{采芙蓉}
\textbf{太宗皇帝}\par\vspace{0.5em}
結伴戲方塘，攜手上雕航。\par
船移分細浪，風散動浮香。\par
遊鸎無定曲，驚鳧有亂行。\par
蓮稀釧聲斷，水廣櫂歌長。\par
棲烏還密樹，泛流歸建章。\par

\section{賦得櫻桃}
\textbf{太宗皇帝}\par\vspace{0.5em}
華林滿芳景，洛陽徧陽春。\par
朱顏含遠日，翠色影長津。\par
喬柯囀嬌鳥，低枝映美人。\par
昔作園中實，今來席上珍。\par

\section{賦得李}
\textbf{太宗皇帝}\par\vspace{0.5em}
玉衡流桂圃，成蹊正可尋。\par
鸎啼密葉外，蝶戲脆花心。\par
麗景光朝彩，輕霞散夕陰。\par
暫顧暉章側，還眺靈山林。\par

\section{賦得浮橋}
\textbf{太宗皇帝}\par\vspace{0.5em}
岸曲非千里，橋斜異七星。\par
暫低逢輦度，還高值浪驚。\par
水搖文鷁動，纜轉錦花縈。\par
遠近隨輪影，輕重應人行。\par

\section{謁幷州大興國寺詩}
\textbf{太宗皇帝}\par\vspace{0.5em}
回鑾遊福地，極目玩芳晨。\par
梵鐘交二響，法日轉雙輪。\par
寶剎遙承露，天花近足春。\par
未佩蘭猶小，無絲柳尚新。\par
圓光低月殿，碎影亂風筠。\par
對此留餘想，超然離俗塵。\par

\section{詠興國寺佛殿前幡}
\textbf{太宗皇帝}\par\vspace{0.5em}
拂霞疑電落，騰虛狀寫虹。\par
屈伸煙霧裏，低舉白雲中。\par
紛披乍依迥，掣曳或隨風。\par
念茲輕薄質，無翅強搖空。\par

\section{望送魏徵葬}
\textbf{太宗皇帝}\par\vspace{0.5em}
閶闔總金鞍，上林移玉輦。\par
野郊愴新別，河橋非舊餞。\par
慘日映峰沈，愁雲隨蓋轉。\par
哀笳時斷續，悲旌乍舒卷。\par
望望情何極，浪浪淚空泫。\par
無復昔時人，芳春共誰遣。\par

\section{傷遼東戰亡}
\textbf{太宗皇帝}\par\vspace{0.5em}
鑿門初奉律，仗戰始臨戎。\par
振鱗方躍浪，騁翼正凌風。\par
未展六奇術，先虧一簣功。\par
防身豈乏智，殉命有餘忠。\par

\section{月晦}
\textbf{太宗皇帝}\par\vspace{0.5em}
晦魄移中律，凝暄起麗城。\par
罩雲朝蓋上，穿露曉珠呈。\par
笑樹花分色，啼枝鳥合聲。\par
披襟歡眺望，極目暢春情。\par

\section{秋日翠微宮}
\textbf{太宗皇帝}\par\vspace{0.5em}
秋光凝翠嶺，涼吹肅離宮。\par
荷疎一蓋缺，樹冷半帷空。\par
側陣移鴻影，圓花釘菊叢。\par
攄懷俗塵外，高眺白雲中。\par

\section{初秋夜坐}
\textbf{太宗皇帝}\par\vspace{0.5em}
斜廊連綺閣，初月照宵幃。\par
塞冷鴻飛疾，園秋蟬噪遲。\par
露結林疎葉，寒輕菊吐滋。\par
愁心逢此節，長歎獨含悲。\par

\section{秋日二首 一}
\textbf{太宗皇帝}\par\vspace{0.5em}
菊散金風起，荷疎玉露圓。\par
將秋數行雁，離夏幾林蟬。\par
雲凝愁半嶺，霞碎纈高天。\par
還似成都望，直見峨眉前。\par

\section{秋日二首 二}
\textbf{太宗皇帝}\par\vspace{0.5em}
爽氣澄蘭沼，秋風動桂林。\par
露凝千片玉，菊散一叢金。\par
日岫高低影，雲空點綴陰。\par
蓬瀛不可望，泉石且娛心。\par

\section{冬宵各爲四韻}
\textbf{太宗皇帝}\par\vspace{0.5em}
雕宮靜龍漏，綺閣宴公侯。\par
珠簾燭燄動，繡柱月光浮。\par
雲起將歌發，風停與管遒。\par

\section{冬日臨昆明池}
\textbf{太宗皇帝}\par\vspace{0.5em}
石鯨分玉溜，劫燼隱平沙。\par
柳影冰無葉，梅心凍有花。\par
寒野凝朝霧，霜天散夕霞。\par
歡情猶未極，落景遽西斜。\par

\section{望雪}
\textbf{太宗皇帝}\par\vspace{0.5em}
凍雲宵徧嶺，素雪曉凝華。\par
入牖千重碎，迎風一半斜。\par
不妝空散粉，無樹獨飄花。\par
縈空慙夕照，破彩謝晨霞。\par

\section{守歲}
\textbf{太宗皇帝}\par\vspace{0.5em}
暮景斜芳殿，年華麗綺宮。\par
寒辭去冬雪，暖帶入春風。\par
階馥舒梅素，盤花卷燭紅。\par
共歡新故歲，迎送一宵中。\par

\section{除夜}
\textbf{太宗皇帝}\par\vspace{0.5em}
歲陰窮暮紀，獻節啓新芳。\par
冬盡今宵促，年開明日長。\par
冰消出鏡水，梅散入風香。\par
對此歡終宴，傾壺待曙光。\par

\section{詠雨}
\textbf{太宗皇帝}\par\vspace{0.5em}
和氣吹綠野，梅雨灑芳田。\par
新流添舊澗，宿霧足朝煙。\par
雁溼行無次，花霑色更鮮。\par
對此欣登歲，披襟弄五弦。\par

\section{賦得含峰雲}
\textbf{太宗皇帝}\par\vspace{0.5em}
翠樓含曉霧，蓮峰帶晚雲。\par
玉葉依巖聚，金枝觸石分。\par
橫天結陣影，逐吹起羅文。\par
非復陽臺下，空將惑楚君。\par

\section{三層閣上置音聲}
\textbf{太宗皇帝}\par\vspace{0.5em}
綺筵移暮景，紫閣引宵煙。\par
隔棟歌塵合，分階舞影連。\par
聲流三處管，響亂一重絃。\par
不似秦樓上，吹簫空學仙。\par

\section{遠山澄碧霧}
\textbf{太宗皇帝}\par\vspace{0.5em}
殘雲收翠嶺，夕霧結長空。\par
帶岫凝全碧，障霞隱半紅。\par
髣髴分初月，飄颻度曉風。\par
還因三里處，冠蓋遠相通。\par

\section{賦得花庭霧}
\textbf{太宗皇帝}\par\vspace{0.5em}
蘭氣已熏宮，新橤半粧叢。\par
色含輕重霧，香引去來風。\par
拂樹濃舒碧，縈花薄蔽紅。\par
還當雜行雨，髣髴隱遙空。\par

\section{春池柳}
\textbf{太宗皇帝}\par\vspace{0.5em}
年柳變池臺，隋堤曲直回。\par
逐浪絲陰去，迎風帶影來。\par
疎黃一鳥弄，半翠幾眉開。\par
縈雪臨春岸，參差間早梅。\par

\section{芳蘭}
\textbf{太宗皇帝}\par\vspace{0.5em}
春暉開紫苑，淑景媚蘭場。\par
映庭含淺色，凝露泫浮光。\par
日麗參差影，風傳輕重香。\par
會須君子折，佩裏作芬芳。\par

\section{詠桃}
\textbf{太宗皇帝}\par\vspace{0.5em}
禁苑春暉麗，花蹊綺樹妝。\par
綴條深淺色，點露參差光。\par
向日分千笑，迎風共一香。\par
如何仙嶺側，獨秀隱遙芳。\par

\section{賦簾}
\textbf{太宗皇帝}\par\vspace{0.5em}
參差垂玉闕，舒卷映蘭宮。\par
珠光搖素月，竹影亂清風。\par
彩散銀鉤上，文斜桂戶中。\par
惟當雜羅綺，相與媚房櫳。\par

\section{詠烏代陳師道}
\textbf{太宗皇帝}\par\vspace{0.5em}
凌晨麗城去，薄暮上林棲。\par
辭枝枝暫起，停樹樹還低。\par
向日終難託，迎風詎肯迷。\par
只待纖纖手，曲裏作宵啼。\par

\section{詠飲馬}
\textbf{太宗皇帝}\par\vspace{0.5em}
駿骨飲長涇，奔流灑絡纓。\par
細紋連噴聚，亂荇繞蹄縈。\par
水光鞍上側，馬影溜中橫。\par
翻似天池裏，騰波龍種生。\par

\section{賦得殘菊}
\textbf{太宗皇帝}\par\vspace{0.5em}
階蘭凝曙霜，岸菊照晨光。\par
露濃晞晚笑，風勁淺殘香。\par
細葉凋輕翠，圓花飛碎黃。\par
還持今歲色，復結後年芳。\par

\end{document}